\chapter{Specimen Glossary}
This chapter applies the forward transliterator to the opening of the typed passage. Each entry has three lines: the Arabic-script form with the default settings, a Latin respelling generated from the Arabic form by the guide tool, and the English original. The respelling is a reading aid for the author and for readers who do not read Arabic script. It is not part of the notation.

\begin{quote}\noindent \ar{ذَا رِبْرِشَنْ أُفْ لَاكْ أَنْدْ إِتْسْ كَلْتْشَرَلْ سِمْبْتُمْزْ}\\[1pt]
\textit{dhu reepreshun uv lak und its kulcherul simptumz}\\[1pt]
The repression of lack and its cultural symptoms.
\end{quote}
\begin{quote}\noindent \ar{إِنْ ذِسْ بَارْتْ أُفْ هِيلِنْ رُولِنْزْ سَايْكُوسِينِمَا}\\[1pt]
\textit{in dhis paart uv helun raalinz saykohsinumu}\\[1pt]
In this part of Helen Rollins' psychocinema,
\end{quote}
\begin{quote}\noindent \ar{شِي دِلْفْزْ إِنْتُو ذَا نَيْتْشَر أُفْ لَاكْ}\\[1pt]
\textit{shee delvz intoo dhu neycher uv lak}\\[1pt]
she delves into the nature of lack.
\end{quote}

\section{Ambiguity and context}
Some distinctions disappear in the notation. \emph{Is} and \emph{as} share a written form, and so do \emph{night} and \emph{knight}. Taken one at a time, such collisions show where the system's limits lie, and Chapter~\ref{ch:reader} measures them. Their practical weight is a separate question, because readers rarely meet words in isolation. In a sentence, the intended word is usually clear from its neighbours, as it is for homophones in spoken English, and a collision matters only where it also causes confusion in continuous prose. The lists below are therefore best read as evidence about the notation's resolution, and not as a prediction that each pair will cause trouble in use.

\section{The reference vocabulary}
The tables below group words by the feature they illustrate. Every form is generated by the toolkit with the default options, and every table is regenerated whenever the rules change, so the book cannot drift from the code.

\subsection{Doubled o and related spellings}
\begin{tabular}{@{}lll@{}}
\toprule
English & Arabic script & Respelling\\
\midrule
book & \ar{بُوكْ} & \textit{buk}\\
look & \ar{لُوكْ} & \textit{luk}\\
good & \ar{غُودْ} & \textit{gud}\\
foot & \ar{فُوتْ} & \textit{fut}\\
wood & \ar{وُودْ} & \textit{wud}\\
food & \ar{فُودْ} & \textit{food}\\
moon & \ar{مُونْ} & \textit{moon}\\
school & \ar{سْكُولْ} & \textit{skool}\\
too & \ar{تُو} & \textit{too}\\
two & \ar{تُو} & \textit{too}\\
blue & \ar{بْلُو} & \textit{bloo}\\
true & \ar{تْرُو} & \textit{troo}\\
\bottomrule
\end{tabular}

\subsection{ee and ea}
\begin{tabular}{@{}lll@{}}
\toprule
English & Arabic script & Respelling\\
\midrule
see & \ar{سِي} & \textit{see}\\
tree & \ar{تْرِي} & \textit{tree}\\
green & \ar{غْرِينْ} & \textit{green}\\
sleep & \ar{سْلِيبْ} & \textit{sleep}\\
head & \ar{هِدْ} & \textit{hed}\\
bread & \ar{بْرِدْ} & \textit{bred}\\
read & \ar{رِدْ} & \textit{red}\\
meet & \ar{مِيتْ} & \textit{meet}\\
meat & \ar{مِيتْ} & \textit{meet}\\
dream & \ar{دْرِيمْ} & \textit{dreem}\\
dead & \ar{دِدْ} & \textit{ded}\\
\bottomrule
\end{tabular}

\subsection{Glides and diphthongs}
\begin{tabular}{@{}lll@{}}
\toprule
English & Arabic script & Respelling\\
\midrule
rain & \ar{رَيْنْ} & \textit{reyn}\\
day & \ar{دَيْ} & \textit{dey}\\
light & \ar{لَايْتْ} & \textit{layt}\\
night & \ar{نَايْتْ} & \textit{nayt}\\
voice & \ar{فُويْسْ} & \textit{voys}\\
flower & \ar{فْلَاوَرْ} & \textit{flower}\\
hour & \ar{آوَرْ} & \textit{ower}\\
boy & \ar{بُويْ} & \textit{boy}\\
toy & \ar{تُويْ} & \textit{toy}\\
noise & \ar{نُويْزْ} & \textit{noyz}\\
now & \ar{نَاوْ} & \textit{now}\\
how & \ar{هَاوْ} & \textit{how}\\
\bottomrule
\end{tabular}

\subsection{Short vowels}
\begin{tabular}{@{}lll@{}}
\toprule
English & Arabic script & Respelling\\
\midrule
cat & \ar{كَاتْ} & \textit{kat}\\
bed & \ar{بِدْ} & \textit{bed}\\
sit & \ar{سِتْ} & \textit{sit}\\
hot & \ar{هُوتْ} & \textit{haat}\\
cut & \ar{كَتْ} & \textit{kut}\\
put & \ar{بُتْ} & \textit{put}\\
pen & \ar{بِنْ} & \textit{pen}\\
map & \ar{مَابْ} & \textit{map}\\
top & \ar{تُوبْ} & \textit{taap}\\
sun & \ar{سَنْ} & \textit{sun}\\
\bottomrule
\end{tabular}

\subsection{Silent e}
\begin{tabular}{@{}lll@{}}
\toprule
English & Arabic script & Respelling\\
\midrule
make & \ar{مَيْكْ} & \textit{meyk}\\
time & \ar{تَايْمْ} & \textit{taym}\\
home & \ar{هُومْ} & \textit{hohm}\\
use & \ar{يُوسْ} & \textit{yoos}\\
cake & \ar{كَيْكْ} & \textit{keyk}\\
like & \ar{لَايْكْ} & \textit{layk}\\
phone & \ar{فُونْ} & \textit{fohn}\\
tune & \ar{تُونْ} & \textit{toon}\\
\bottomrule
\end{tabular}

\subsection{R-coloured vowels}
\begin{tabular}{@{}lll@{}}
\toprule
English & Arabic script & Respelling\\
\midrule
car & \ar{كَارْ} & \textit{kaar}\\
far & \ar{فَارْ} & \textit{faar}\\
her & \ar{هِرْ} & \textit{her}\\
bird & \ar{بِرْدْ} & \textit{berd}\\
turn & \ar{تَرْنْ} & \textit{tern}\\
word & \ar{وُرْدْ} & \textit{werd}\\
door & \ar{دُورْ} & \textit{dawr}\\
more & \ar{مُورْ} & \textit{mawr}\\
fire & \ar{فَايَرْ} & \textit{fayer}\\
tire & \ar{تَايَرْ} & \textit{tayer}\\
\bottomrule
\end{tabular}

\subsection{Consonants}
\begin{tabular}{@{}lll@{}}
\toprule
English & Arabic script & Respelling\\
\midrule
thing & \ar{ثِنْغْ} & \textit{thing}\\
think & \ar{ثِنْغْكْ} & \textit{thingk}\\
this & \ar{ذِسْ} & \textit{dhis}\\
that & \ar{ذَتْ} & \textit{dhat}\\
church & \ar{تْشَرْتْشْ} & \textit{cherch}\\
change & \ar{تْشَيْنْجْ} & \textit{cheynj}\\
judge & \ar{جَجْ} & \textit{juj}\\
measure & \ar{مِيزَرْ} & \textit{mezher}\\
vision & \ar{فِيزَنْ} & \textit{vizhun}\\
singer & \ar{سِنْغَرْ} & \textit{singer}\\
finger & \ar{فِنْغَرْ} & \textit{fingger}\\
garden & \ar{غَارْدِنْ} & \textit{gaardun}\\
\bottomrule
\end{tabular}

\subsection{Clusters}
\begin{tabular}{@{}lll@{}}
\toprule
English & Arabic script & Respelling\\
\midrule
string & \ar{سْتْرِنْغْ} & \textit{string}\\
spring & \ar{سْبْرِنْغْ} & \textit{spring}\\
splash & \ar{سْبْلَاشْ} & \textit{splash}\\
street & \ar{سْتْرِيتْ} & \textit{street}\\
strength & \ar{سْتْرِنْغْكْثْ} & \textit{strengkth}\\
twelfth & \ar{تْوِلْفْثْ} & \textit{twelfth}\\
asked & \ar{أَسْكْتْ} & \textit{askt}\\
texts & \ar{تِكْسْتْسْ} & \textit{teksts}\\
\bottomrule
\end{tabular}

\section{Teaching sequence}
The materials in this chapter suit a staged course in which the supporting lines are withdrawn as the reader gains fluency.
\begin{enumerate}
\item \textbf{Paired reading.} Begin with the Arabic-script English beside the familiar English word, so that each new spelling is tied to a word the reader already knows.
\item \textbf{A small recurring vocabulary.} Practise a few dozen frequent words (\emph{the}, \emph{of}, \emph{and}, \emph{is}, \emph{in}, \emph{to}) until their written forms are recognised at once. Their fixed spellings make this easy.
\item \textbf{Short passages with the English covered.} Read short passages with the English line concealed, uncovering it only to check.
\item \textbf{Withdraw the respelling.} Drop the Latin respelling next, since it is an approximate aid. It introduces spelling conventions of its own (the \emph{ee} of \emph{reeding}, the \emph{u} of \emph{uv}), and a reader who leans on it too long learns those conventions instead of the notation.
\item \textbf{Lighter marks.} Finally, move from full marks to medial, vowel-only and unmarked text, using the density settings.
\end{enumerate}

\section{Reading practice}
The following passages start with the simplest sentences and become longer. Readers learning the notation can cover the English and the Latin respelling and read the Arabic-script line aloud.

\begin{quote}\noindent \ar{ذَا بُوكْ إِزْ رِتِنْ إِنْ إِنْغْلِشْ.}\\[2pt]
\ar{ذا بوك إز رتن إن إنغلش.}\\[1pt]
\textit{dhu buk iz ritun in ingglish}\\[1pt]
The book is written in English.
\end{quote}
\begin{quote}\noindent \ar{ذَا لِتَرْزْ تْشَيْنْجْ، بَتْ ذَا وُرْدْزْ رِمَيْنْ.}\\[2pt]
\ar{ذا لترز تشينج، بت ذا وردز رمين.}\\[1pt]
\textit{dhu leterz cheynj but dhu werdz rimeyn}\\[1pt]
The letters change, but the words remain.
\end{quote}
\begin{quote}\noindent \ar{أَ رِيدَرْ لَرْنْزْ ذَا سْكْرِبْتْ أَنْدْ بِغِنْزْ تُو رِيكُغْنَايْزْ أَ فَمِيلْيَرْ فُويْسْ.}\\[2pt]
\ar{أ ريدر لرنز ذا سكربت أند بغنز تو ريكغنايز أ فميلير فويس.}\\[1pt]
\textit{u reeder lernz dhu skript und biginz too rekugnayz u fumilyer voys}\\[1pt]
A reader learns the script and begins to recognise a familiar voice.
\end{quote}
\begin{quote}\noindent \ar{أَلِسْ وَازْ بِغِنِنْغْ تُو غِتْ فِيرِي تَايَرْدْ أُفْ سِتِنْغْ بَايْ هِرْ سِسْتَرْ أُونْ ذَا بَانْغْكْ، أَنْدْ أُفْ هَافِنْغْ نُوثِنْغْ تُو دُو.}\\[2pt]
\ar{ألس واز بغننغ تو غت فيري تايرد أف ستنغ باي هر سستر أون ذا بانغك، أند أف هافنغ نوثنغ تو دو.}\\[1pt]
\textit{alus waaz bigining too get veree tayerd uv siting bay her sister aan dhu bangk und uv having nuthing too doo}\\[1pt]
Alice was beginning to get very tired of sitting by her sister on the bank, and of having nothing to do.
\end{quote}
\begin{quote}\noindent \ar{وُونْسْ أُورْ تْوَايْسْ شِي هَادْ بِيبْتْ إِنْتُو ذَا بُوكْ هِرْ سِسْتَرْ وَازْ رِيدِنْغْ، بَتْ إِتْ هَادْ نُو بِكْتْشَرْزْ أُورْ كُونْفَرْسَيْشَنْزْ إِنْ إِتْ.}\\[2pt]
\ar{وونس أور توايس شي هاد بيبت إنتو ذا بوك هر سستر واز ريدنغ، بت إت هاد نو بكتشرز أور كونفرسيشنز إن إت.}\\[1pt]
\textit{wuns awr tways shee had peept intoo dhu buk her sister waaz reeding but it had noh pikcherz awr kaanverseyshunz in it}\\[1pt]
Once or twice she had peeped into the book her sister was reading, but it had no pictures or conversations in it.
\end{quote}
\begin{quote}\noindent \ar{أَنْدْ وَتْ إِزْ ذَا يُوسْ أُفْ أَ بُوكْ، ثُوتْ أَلِسْ، وِثَاوْتْ بِكْتْشَرْزْ أُورْ كُونْفَرْسَيْشَنْزْ؟}\\[2pt]
\ar{أند وت إز ذا يوس أف أ بوك، ثوت ألس، وثاوت بكتشرز أور كونفرسيشنز؟}\\[1pt]
\textit{und wut iz dhu yoos uv u buk thawt alus withowt pikcherz awr kaanverseyshunz}\\[1pt]
And what is the use of a book, thought Alice, without pictures or conversations?
\end{quote}
\begin{quote}\noindent \ar{إِنْ ذَا بِغِنِنْغْ غُودْ كْرَيْتَدْ ذَا هِيفِنْ أَنْدْ ذَا أَرْثْ.}\\[1pt]
\textit{in dhu bigining gaad kreeeytud dhu hevun und dhu erth}\\[1pt]
In the beginning God created the heaven and the earth.
\end{quote}
\begin{quote}\noindent \ar{أَنْدْ ذَا أَرْثْ وَازْ وِثَاوْتْ فُورْمْ، أَنْدْ فُويْدْ؛ أَنْدْ دَارْكْنِسْ وَازْ أَبُونْ ذَا فَيْسْ أُفْ ذَا دِيبْ.}\\[1pt]
\textit{und dhu erth waaz withowt fawrm und voyd und daarknus waaz upaan dhu feys uv dhu deep}\\[1pt]
And the earth was without form, and void; and darkness was upon the face of the deep.
\end{quote}
\begin{quote}\noindent \ar{أَنْدْ غُودْ سِدْ، لِتْ ذِرْ بِي لَايْتْ: أَنْدْ ذِرْ وَازْ لَايْتْ.}\\[1pt]
\textit{und gaad sed let dher bee layt und dher waaz layt}\\[1pt]
And God said, Let there be light: and there was light.
\end{quote}
\begin{quote}\noindent \ar{ذَا تْرَافِلَرْ كَيْمْ دَاوْنْ فْرُومْ ذَا سْتَارْ سِرَسْ تُو آوَرْ لِتَلْ أَنْتْ هِلْ، أَنْدْ مِزَرْدْ ذَا أَرْثْ وِذْ هِزْ فِنْغَرْزْ.}\\[1pt]
\textit{dhu travuler keym down frum dhu staar sirius too ower litul ant hil und mezherd dhu erth widh hiz finggerz}\\[1pt]
The traveller came down from the star Sirius to our little ant hill, and measured the earth with his fingers.
\end{quote}
\begin{quote}\noindent \ar{هِي فَاوْنْدْ ذَا بِيبَلْ سْمُولْ، بَتْ هِي فَاوْنْدْ ذِمْ كْيُرِيَسْ، أَنْدْ هِي أَسْكْتْ ذِمْ مِنِي كْوِسْتْشَنْزْ أَبَاوْتْ ذِرْ لِفْزْ.}\\[1pt]
\textit{hee fownd dhu peepul smawl but hee fownd dhem kyureeus und hee askt dhem menee kweschunz ubowt dher livz}\\[1pt]
He found the people small, but he found them curious, and he asked them many questions about their lives.
\end{quote}
